% **************************************************
% * Opciones de clase y globales
% **************************************************
\documentclass[%
    a4paper,               % Tamaño de papel A4
    twoside,               % Páginas pares e impares con diseño diferenciado
    10pt,                  % Tamaño de fuente. No inferior a 10 puntos.
]{article}                 % Basado en el tipo "article"

% **************************************************
% * Idioma y paquete de formato
% **************************************************
\usepackage[galician]{babel}
\usepackage{tfgteleco}

% **************************************************
% * Datos para la portada del TFG
% **************************************************
\title{Deseño e desenvolvemento dunha funcionalidade de busca e filtrado de cursos no catálogo de cursos en liña}
\author{Aarón Riveiro Vilar}
\titores{Manuel Caeiro Rodríguez}
\curso{2025/2026}

% **************************************************
% * Cabeceras (opcionales)
% **************************************************
\cabeceiras{true}

% **************************************************
% * Inicio del documento
% **************************************************
\begin{document}

\hypersetup{pageanchor=false}
\maketitle
% **************************************************
% * Tábla de contenidos (opcional)
% **************************************************
\tableofcontents
\clearpage
\hypersetup{pageanchor=true}

% **************************************************
% * Introducción
% **************************************************
\section{Introdución}
Nos últimos anos, o crecemento sostido dos sistemas de información dixitais deu lugar a repositorios con volumes de datos cada vez maiores e máis heteroxéneos. Neste contexto, a mera dispoñibilidade da información xa non resulta suficiente: é necesario dotar os sistemas de mecanismos eficaces que permitan localizar, filtrar e presentar os datos de forma relevante para o usuario. A disciplina da recuperación de información (\textit{Information Retrieval}) estuda precisamente os métodos e modelos que permiten acceder de maneira eficiente a grandes coleccións de documentos, maximizando a relevancia dos resultados ofrecidos fronte a unha consulta determinada \cite{information_retrieval}.

A medida que aumenta o tamaño e a complexidade dos repositorios, a experiencia de usuario pasa a depender en grande medida da calidade do sistema de busca e filtrado. Non só é importante a rapidez de resposta, senón tamén a precisión e exhaustividade dos resultados, así como a capacidade do sistema para interpretar a intención de busca do usuario e ofrecer mecanismos de refinamento progresivo mediante filtros, facetas ou criterios adicionais \cite{retrieval_accuracy}. Un sistema de busca limitado a coincidencias literais sobre un único campo de texto pode resultar insuficiente en contornas nas que os datos presentan múltiples dimensións e atributos relevantes.

Por iso, o deseño e a implementación de estratexias avanzadas de busca e filtrado constitúen un elemento clave en calquera sistema de información orientado ao usuario final, especialmente cando se pretende facilitar o acceso eficiente a contidos estruturados e mellorar a usabilidade global da plataforma.

% **************************************************
% * Contexto
% **************************************************
\subsection{Contexto do proxecto}
O proxecto DACEM (\textit{Digitising Academic Catalogues for Enhanced Mobility}) \cite{dacem}, aprobado na convocatoria de 2023 do programa Erasmus+ para Asociacións Estratéxicas en Educación Superior (código de proxecto 2023-1-ES01-KA220-HED-000160344), ten como finalidade o desenvolvemento dunha plataforma de catálogo de cursos dirixida a institucións educativas europeas. O seu obxectivo principal é ofrecer un sistema que permita publicar e manter información actualizada sobre titulacións e materias de forma accesible, estruturada e interoperable, facilitando así os procesos de mobilidade académica internacional. A plataforma está orientada a dous perfís principais de usuarios: as institucións educativas, que a empregan para publicar e manter actualizado o seu catálogo de cursos; e os estudantes que, no marco de procesos de mobilidade Erasmus, consultan a oferta formativa das universidades parceiras.

O sistema de información asociado a esta plataforma está fundamentado en tecnoloxías de software libre e está concibido para poder despregarse tanto en contornas locais como en infraestruturas na nube. Máis alá do almacenamento estruturado de datos académicos, o correcto funcionamento do sistema depende da súa capacidade para xestionar, organizar e recuperar a información de maneira eficiente. Neste contexto, os mecanismos de busca e filtrado desempeñan un papel esencial, xa que constitúen o principal medio de interacción entre o usuario e o repositorio de cursos. A calidade destes mecanismos condiciona directamente a usabilidade do sistema, a precisión dos resultados obtidos e a experiencia global de navegación.

% **************************************************
% * Motivación
% **************************************************
\subsection{Motivación}
Un dos retos fundamentais nos sistemas de información que xestionan grandes volumes de datos estruturados non reside unicamente no seu almacenamento, senón na forma en que estes datos son recuperados e presentados ao usuario. A medida que un repositorio medra en número de elementos e en complexidade de atributos, os mecanismos de busca simples baseados en coincidencias textuais directas poden resultar insuficientes para ofrecer resultados relevantes e axustados ás necesidades reais de consulta.

Esta problemática é especialmente significativa en contornas nas que os datos presentan múltiples dimensións descritivas —como titulacións, materias, áreas de coñecemento, idiomas, créditos ou modalidades— e nas que o usuario pode ter criterios de busca diversos e combinados. A ausencia de mecanismos avanzados de filtrado e refinamento progresivo limita a capacidade do sistema para adaptarse a diferentes perfís de usuario e escenarios de consulta, reducindo a eficiencia na localización de información pertinente.

No ámbito actual do desenvolvemento de aplicacións web e sistemas de xestión de contidos, as tendencias apuntan cara á implementación de motores de busca máis sofisticados, capaces de indexar múltiples campos, ponderar resultados segundo a relevancia, incorporar criterios de filtrado dinámico e mellorar a experiencia de usuario mediante interfaces intuitivas. Estas solucións non só optimizan o acceso á información, senón que contribúen a unha maior usabilidade e aproveitamento do sistema no seu conxunto.

O presente Traballo de Fin de Grao enmárcase neste contexto, propoñendo o deseño e desenvolvemento dunha mellora nos mecanismos de busca e filtrado do catálogo de cursos da plataforma DACEM. O obxectivo é evolucionar dende un sistema de coincidencia textual básica cara a unha solución máis flexible e precisa, aliñada coas boas prácticas actuais en recuperación de información e desenvolvemento web, empregando tecnoloxías abertas e integradas na contorna existente. Concretamente, abórdase a implementación de criterios de ordenación de resultados para as vistas de busca existentes, así como un mecanismo de busca difusa que permita localizar resultados aínda en presenza de erros tipográficos na consulta.

% **************************************************
% * Plataforma de partida
% **************************************************
\subsection{A plataforma de partida}

A plataforma DACEM está implementada sobre Drupal 10, un sistema de xestión de contidos (CMS) de código aberto amplamente empregado no ámbito académico e institucional. Drupal 10 proporciona unha contorna flexible para a definición de tipos de contido estruturado, a xestión de usuarios e a publicación de información a través de interfaces web configurables.

O modelo de datos organízase arredor de tres tipos de contido principais: \textit{programme}, que representa unha titulación académica; \textit{individual educational component} (IEC), correspondente a unha materia ou compoñente curricular individual; e \textit{institution}, que modela as universidades participantes. Cada entidade conta cun conxunto de campos estruturados que recollen atributos como a área de coñecemento, o nivel de cualificación, os créditos, o idioma de impartición ou o país de orixe, entre outros.

Para a presentación e consulta do catálogo, a plataforma emprega o módulo \textit{Views}, integrado no núcleo de Drupal, que permite definir listaxes e consultas de contido de forma declarativa. O módulo \textit{Better Exposed Filters} complementa \textit{Views} proporcionando a interface de filtrado ao usuario. Existen dúas vistas principais de busca: \textit{search\_programme}, accesible en \texttt{/search-programme}, orientada á consulta de titulacións; e \textit{search\_iec}, en \texttt{/search-iec}, orientada á consulta de compoñentes educativos individuais. A capa de presentación visual está baseada nun tema personalizado construído sobre Bootstrap 5.

O mecanismo de busca textual actual baséase no filtro \textit{combine} de \textit{Views}, que realiza comparacións de subcadeas sobre o campo título das entidades. Esta aproximación presenta limitacións relevantes: a busca restrinxese ao campo título, ignorando atributos como a descrición ou a área de coñecemento; non se aplica ningunha ponderación por relevancia, polo que os resultados se presentan en orde alfabética con independencia da pertinencia da consulta; e non existe tolerancia a erros tipográficos, de xeito que calquera inexactitude na consulta pode producir cero resultados aínda que o elemento buscado exista no catálogo. Ademais, as opcións de filtrado son estáticas e non se adaptan ao conxunto de resultados dispoñibles, polo que poden aparecer valores sen correspondencia con ningunha entrada existente.

% **************************************************
% * Objetivos y requisitos
% **************************************************
\section{Obxectivos e requisitos}
O presente traballo céntrase na mellora dos mecanismos de busca e filtrado do catálogo de cursos desenvolvido no marco do proxecto DACEM. A diferenza doutros desenvolvementos orientados á recompilación ou estruturación de datos, este proxecto aborda a fase de recuperación e presentación da información, poñendo o foco na interacción entre o usuario e o sistema.

O obxectivo principal consiste en analizar as limitacións do sistema actual de consulta e implementar unha solución que permita mellorar a precisión, flexibilidade e usabilidade dos resultados obtidos. Para iso, defínense a continuación os requisitos funcionais e non funcionais que deben guiar o deseño e a implementación da proposta.

Enténdese por requisito funcional aquel que describe as capacidades concretas que debe ofrecer o sistema, mentres que os requisitos non funcionais establecen condicións relativas ao seu comportamento, rendemento, integración e manexabilidade, sen definir directamente novas funcionalidades.

% **************************************************
% * Requisitos funcionales
% **************************************************
\subsection{Requisitos funcionais}

% **************************************************
% * Sistema de búsqueda
% **************************************************
\subsubsection{Sistema de busca}
\begin{enumerate}
    \item O sistema debe permitir a realización de buscas textuais sobre o catálogo de cursos a partir dunha cadea de consulta introducida polo usuario.

    \item A busca non debe limitarse exclusivamente ao título do curso, senón que deberá contemplar múltiples campos relevantes do modelo de datos, tales como a descrición, a área de coñecemento, a titulación asociada ou outros atributos significativos que formen parte da información estruturada do catálogo.

    \item O sistema debe soportar coincidencias parciais e non estritamente literais, permitindo a localización de resultados mesmo cando a consulta non coincida exactamente co texto almacenado. Neste sentido, deberase contemplar a normalización de maiúsculas e minúsculas e, no seu caso, a xestión de caracteres acentuados.

    \item O sistema deberá ordenar os resultados da busca segundo un criterio de relevancia, de maneira que os elementos máis pertinentes en relación coa consulta aparezan nas primeiras posicións.

    \item O sistema debe garantir que as buscas con consultas baleiras ou ambiguas non provoquen erros nin comportamentos inesperados, definindo un comportamento consistente nestes casos (por exemplo, amosando todos os resultados ou solicitando unha consulta válida).

          % // \item O sistema deberá integrarse co modelo de datos existente sen alterar a súa estrutura fundamental, reutilizando na medida do posible os mecanismos de indexación e consulta proporcionados pola plataforma.

          % // \item O sistema debe permitir a extensión futura do mecanismo de busca, de maneira que poidan incorporarse novos campos ou criterios sen requirir unha reestruturación completa da solución implementada.
\end{enumerate}

% **************************************************
% * Sistema de filtrado
% **************************************************
\subsubsection{Sistema de filtrado}
\begin{enumerate}
    \item O sistema debe permitir a aplicación de filtros sobre o conxunto de resultados obtidos, posibilitando a restrición dos cursos amosados en función de distintos atributos do modelo de datos.

    \item Os filtros deberán contemplar, cando menos, aqueles campos estruturados que resulten relevantes para a consulta dos usuarios, tales como titulación, área de coñecemento, idioma de impartición, número de créditos, modalidade ou outros atributos dispoñibles no catálogo.

    \item O sistema debe permitir a combinación simultánea de múltiples filtros, de maneira que o usuario poida refinar progresivamente os resultados mediante criterios acumulativos.

    \item A aplicación ou modificación de filtros non deberá xerar inconsistencias nos resultados amosados, garantindo que o conxunto final de cursos cumpra todos os criterios activos en cada momento.

          % // \item O sistema debe actualizar os resultados de forma coherente tras a activación ou desactivación de filtros, mantendo unha experiencia de usuario fluída e evitando recargas innecesarias que afecten negativamente á interacción.

    \item O sistema debe permitir que os filtros sexan dinámicos e adaptativos ao conxunto de datos dispoñible, evitando amosar opcións de filtrado que non se correspondan con resultados existentes.

          % // \item O mecanismo de filtrado deberá integrarse de forma consistente co sistema de busca, permitindo que ambos procesos actúen de maneira complementaria, de xeito que a consulta textual e os criterios estruturados poidan aplicarse conxuntamente.

          % // \item O deseño do sistema de filtrado debe contemplar a posibilidade de ampliación futura, permitindo a incorporación de novos criterios sen necesidade de redeseñar a arquitectura principal.
\end{enumerate}

% **************************************************
% * Requisitos no funcionales
% **************************************************
\subsection{Requisitos non funcionais}

% **************************************************
% * Integración y rendimiento
% **************************************************
\subsubsection{Integración e rendemento}
\begin{enumerate}
    \item A solución proposta debe integrarse coa plataforma existente sen alterar a estrutura fundamental do modelo de datos nin comprometer o funcionamento do resto de funcionalidades do sistema.

    \item O sistema debe ser compatible coa arquitectura actual baseada en Drupal e cos módulos empregados para a xestión do contido, garantindo que a mellora do mecanismo de busca e filtrado non requira unha reimplementación completa da contorna.

    \item A implementación debe reutilizar, na medida do posible, os mecanismos de indexación e consulta dispoñibles na plataforma, ou ben estendelos de forma controlada cando resulte necesario.

    \item O sistema debe garantir tempos de resposta axeitados nas consultas realizadas polo usuario, de maneira que a experiencia de uso non se vexa degradada por demoras excesivas na xeración de resultados.

    \item O rendemento do sistema debe manterse estable ante un incremento razoable do volume de datos almacenados no catálogo, asegurando que a escalabilidade da solución permita a súa evolución futura.

    \item O sistema debe soportar a execución concorrente de múltiples consultas sen producir erros, bloqueos ou degradacións significativas do rendemento.

    \item A solución implementada deberá minimizar o impacto nos recursos do servidor, evitando sobrecargas innecesarias e optimizando as consultas realizadas á base de datos.
\end{enumerate}

% **************************************************
% * Otros
% **************************************************
\subsubsection{Outros}
\begin{enumerate}
    \item A solución debe estar fundamentada en tecnoloxías de uso libre e código aberto, mantendo coherencia cos principios e coa arquitectura xeral do proxecto DACEM.

    \item O desenvolvemento debe garantir a mantenibilidade do código, estruturando a implementación de forma modular e documentada para facilitar futuras modificacións ou ampliacións.

    \item A solución debe ser compatible con futuras actualizacións da plataforma e dos módulos empregados, evitando dependencias que comprometan a evolución do sistema.

    \item O sistema debe garantir un comportamento estable e previsible ante entradas incorrectas, consultas mal formadas ou combinacións de filtros non válidas, evitando erros visibles para o usuario final.

    \item A interface de busca e filtrado debe cumprir criterios básicos de usabilidade, asegurando que os mecanismos de interacción resulten intuitivos e coherentes co deseño xeral da plataforma.

    \item A solución debe preservar a integridade dos datos existentes, garantindo que as melloras introducidas na capa de consulta non alteren a información almacenada nin a súa consistencia.

    \item O deseño do sistema deberá contemplar a súa posible reutilización ou adaptación noutros contextos similares, favorecendo a portabilidade e a extensibilidade da solución implementada.
\end{enumerate}

% **************************************************
% * Deseño
% **************************************************
\section{Deseño}

% **************************************************
% * Desarrollo
% **************************************************
\section{Desenvolvemento}

% **************************************************
% * Conclusiones
% **************************************************
\section{Conclusións e liñas futuras}

% **************************************************
% * Bibliografía
% **************************************************
\bibliographystyle{IEEEtran}
\bibliography{references}
\clearpage

% **************************************************
% * Anexos
% **************************************************
\appendix
\section{Anexos}

\end{document}